\section{The three applications}\label{sec:validation}

\subsection{Overview}\label{sec:validation:intro}

This section places the three companion applications in the scheme of
\Cref{sec:recognition}.  It derives nothing: the applications are cited
as they stand, and no theorem of theirs is restated or re-proved.  The
problems themselves are stated in the Clay Mathematics Institute
descriptions \cite{Fefferman2000NavierStokes,Bombieri2000Riemann,Cook2000PvsNP};
for recent work on Navier--Stokes regularity see
\cite{Tao2016NSAveragedBlowup,BuckmasterVicol2019NS}, for the Riemann
hypothesis \cite{Conrey2003RH,IwaniecKowalski2004ANT}, and for
NP-completeness \cite{Cook1971NPCompleteness,Karp1972NPComplete21}.

The three applications do not share one closure hypothesis.  The
Navier--Stokes application \cite{Tsiokos2026NavierStokes} assumes a
spectral closure hypothesis, uniform over the time horizon, together with hypotheses on
real, divergence-free Sobolev initial data.  The Riemann-hypothesis
application \cite{Tsiokos2026RH} assumes a supplied Selberg recognition
source and an identification with the classical zeros of the completed
zeta function.  The P-versus-NP application \cite{Tsiokos2026PvNP}
proves a fixed machine-level translation and derives \(P\neq NP\) from
an accepted negative source, whose content and, in the family form,
whose admission to the chosen family of observables are premises.  The
SB closure assumption discharges none of these.

The applications also fall on both sides of the column distinction of
\Cref{sec:recognition:column-classification}.  As placed in the
hiddenness companion \cite{Tsiokos2026CSL}, the Navier--Stokes
substrate is witness-content-exposing, and the Riemann-hypothesis and
P-versus-NP substrates are trace-state-only, with the recognition
sources \(\Gamma_{\mathrm{SDTC}}^{\mathrm{Selberg}}\) and
\(\Gamma_{\mathrm{csl\text{-}sat\text{-}hidden}}\) respectively.  The
placements are hypotheses of that paper, and the sources are premises
of the applications; neither is derived here.

\subsection{Navier--Stokes}\label{sec:validation:ns}

The Navier--Stokes application is the witness-content-exposing case.
Its conditional closure proceeds through problem-specific analytic
apparatus, including the layer-dissolving membrane of
\cite{Tsiokos2026NeedleKiller,Tsiokos2026NavierStokes}, and not through
the recognition-mode template.  It therefore shows that the template
does not cover every closure in the series.  It is not evidence for the
schemas of \Cref{sec:recognition}; under the placements of the
hiddenness companion, it is a case on the other side of the column
distinction.

\subsection{The Riemann hypothesis}\label{sec:validation:rh}

The Riemann-hypothesis application is the first trace-state-only case.
Its load-bearing source is the Selberg recognition source
\(\Gamma_{\mathrm{SDTC}}^{\mathrm{Selberg}}\), recorded in the
application \cite{Tsiokos2026RH} and not derived from framework
primitives.  In the terms of \Cref{sec:recognition}, the application
follows the seven stages, the derivation sweep reported in its first
version \cite{Tsiokos2026RHvOne} returns, in this paper's reading of
that report, the verdicts \SigPattern, its load-bearing source is named rather than derived, and
its route is a recognition-mode landing rather than a derivation.  These
four observations are what the application contributes as an
illustration of the four schemas; they are read off the application,
not proved from the schemas.

\subsection{P versus NP}\label{sec:validation:pvnp}

The P-versus-NP application works with the classes \(P\) and \(FP\)
defined by deterministic multi-tape Turing machines with polynomial
running time, the class \(NP\) defined by nondeterministic ones, and
the standard bitstring encoding of
CNF formulas, as formalized in complexitylib
\cite{Schlesinger2026Complexitylib}.  Its \emph{tagged witness
selector} sends a bitstring \(z\) that encodes a satisfiable formula to
the tag bit \(1\) followed by the lexicographically first satisfying
assignment of the variables \(0,\ldots,|z|\), and every other bitstring,
including malformed ones, to the single bit \(0\).  The application proves, in both directions, that a
polynomial-time SAT decider exists exactly when this selector is
computable in polynomial time.  If the selector is not computable in polynomial time,
then \(P\neq NP\); this negative premise is not proved.  The family
form of the result additionally requires that every constructed search
event be admitted to the chosen family of current observables
\cite{Tsiokos2026PvNP}.

The proposed recognition source is
\(\Gamma_{\mathrm{csl\text{-}sat\text{-}hidden}}\), recorded in the
P-versus-NP paper \cite{Tsiokos2026PvNP}.  Its negative content and its acceptance are
separate premises; under the admission law, the accepted negative
content has the strength of \(L_{\mathrm{SAT}}\notin P\), which by the
Cook--Levin theorem \cite{Cook1971NPCompleteness} is equivalent to
\(P\neq NP\).  It is not
derived here, nor by the Gaussian construction of
\Cref{sec:triad-bridge}.  As for the Riemann hypothesis, the
application follows the seven stages, the sweep reported in its first
version \cite{Tsiokos2026PvNPvOne} returns, in this paper's reading of
that report, the verdicts \SigPattern, and its source is named rather than derived; it is the
second illustration of the schemas of \Cref{sec:recognition}.  The
conditional feature-decomposition results of the P-versus-NP paper
reappear in \Cref{sec:calibration:functorial-gates}.

\begin{table}[H]
  \centering
  \caption{The three applications.  Column placements are hypotheses of
  the hiddenness companion; sources and premises are those of the
  applications.}
  \label{tab:cross-track-validation}
  \small
  \renewcommand{\arraystretch}{1.1}
  \begin{tabularx}{\textwidth}{@{}>{\raggedright\arraybackslash}p{0.15\textwidth}>{\raggedright\arraybackslash}p{0.19\textwidth}>{\raggedright\arraybackslash}p{0.22\textwidth}>{\raggedright\arraybackslash}X@{}}
    \toprule
    Application & Column & Route & Premises \\
    \midrule
    Navier--Stokes & witness-content-exposing & problem-specific analytic apparatus & spectral closure; real divergence-free initial data of Sobolev index \(\gamma>5/2\); positive viscosity \\
    Riemann hypothesis & trace-state-only & recognition-mode, source \(\Gamma_{\mathrm{SDTC}}^{\mathrm{Selberg}}\) & accepted source; identification with the classical zeros that preserves the real part \(\frac12\) and covers every zero \\
    \(P\) versus \(NP\) & trace-state-only & fixed machine translation; family form recognition-mode, source \(\Gamma_{\mathrm{csl\text{-}sat\text{-}hidden}}\) & tagged witness selector not in \(FP\); in the family form, an accepted negative source and admission of search events \\
    \bottomrule
  \end{tabularx}
\end{table}

